\originalchapter{量子Ising链与Majorana费米子}
\label{ch:ising}
\source{8.513第7--15讲组及作业5--7。}
\originalsection{Jordan--Wigner变换}
横场Ising链$H=-J\sum_j\sigma_j^x\sigma_{j+1}^x-h\sum_j\sigma_j^z$。不同站自旋算符对易，费米算符应反对易，因此不能直接令$c_j=\sigma_j^-$. 加上前方自旋串，定义
\begin{equation}
a_j=\left(\prod_{l<j}\sigma_l^z\right)\sigma_j^x,\qquad b_j=\left(\prod_{l<j}\sigma_l^z\right)\sigma_j^y,\qquad c_j=(a_j+\ii b_j)/2.
\end{equation}
对于$j<k$，$a_j$与$k$串中的$\sigma_j^z$反对易，其他因子对易，所以$\{a_j,a_k\}=0$；同站$a_j^2=b_j^2=1$，完整关系$\{a_j,a_k\}=\{b_j,b_k\}=2\delta_{jk}$、$\{a_j,b_k\}=0$。又有$\sigma_j^z=-\ii a_jb_j$、$\sigma_j^x\sigma_{j+1}^x=-\ii b_ja_{j+1}$，因此开链$H=\ii\sum_j(Jb_ja_{j+1}+ha_jb_j)$。

这里$c_j$为非局域自旋算符，虽然费米哈密顿量局域，但自旋局域观测量可能对应费米字符串。这个区别会直接影响散射谱，不能将单费米激发自动认作一个自旋翻转产生的态。
\originalsection{能谱与量子临界点}
对无穷链或固定边界扇区Fourier变换，再作费米Bogoliubov变换，正能谱为
\begin{equation}
E_k=2\sqrt{J^2+h^2-2Jh\cos k}.
\end{equation}
$|h|<|J|$为自旋有序区，$|h|>|J|$为横场极化区，$|h|=|J|$隙闭合。$h=J>0$时$E_k\simeq2J|k|$，动态指数$z=1$，非临界质量$\Delta=2|J-h|$。连续低能是左右运动的实Majorana场。

在开链且$h=0$时，$b_j$只与$a_{j+1}$配对，$a_1$和$b_L$完全不在哈密顿量中，组合成一个零能费米自由度。$0<|h/J|<1$时边缘模$\gamma_L\propto\sum_j(h/J)^{j-1}a_j$，归一要求$|h/J|<1$。两个端模重叠导致$\Delta E\propto\e^{-L/\xi}$，$\xi^{-1}=|\log|h/J||$。零模不是在任意参数下存在的孤立粒子，体隙和边界是必要条件。
\originalsection{临界零模与四次项的上临界维数}
上述一维链的能谱在临界点线性，但把这一Gaussian结果推广到任意空间维数，需要检查被丢弃的相互作用。8.513第7--15讲组第25--28页用$d$维横场Ising的序参量$q$说明这点。在有限体积$N=L^d$、平均场临界位置$h=2dJ$，均匀模的二次$q^2$项消失，低能相空间拉氏量为
\begin{equation}
L_0=Np\dot q-\frac{Nh}{2}p^2-Ngq^4.
\end{equation}
令$\widetilde p=\sqrt Np,\ \widetilde q=\sqrt Nq$，则$H_0=h\widetilde p^2/2+(g/N)\widetilde q^4$。量子不确定性$\widetilde p\widetilde q\sim1$使基态估算为$h/\widetilde q^2+g\widetilde q^4/N$；两项平衡给$\widetilde q\sim N^{1/6}$及$E_0\sim N^{-1/3}$（省去固定的$h,g$系数）。零模无论多高维都不可简单按自由振子处理。

对非零波矢，Gaussian能量$\omega_k\sim J|k|$，其真空涨落$\widetilde q^2\sim1/|k|$。四次校正$\delta E_k\sim g/(Nk^2)$相对谐振能$E_k\sim J|k|$的比值为$\sim g/(JNk^3)$。最低$k\sim L^{-1}=N^{-1/d}$，所以比值随大小如$N^{3/d-1}$：$d>3$趋零，$d=3$临界，$d<3$增长。这个模式估算与连续场论的幂次计数相符。

取动态指数$z=1$，Euclidean作用写$S=\int\dd\tau\,\dd^dx\,[(\partial_\tau\phi)^2+c^2(\nabla\phi)^2+r\phi^2+u\phi^4]/2$，其中四次项系数的常数$1/2$已吸进$u$。保持梯度项无量纲给$[\phi]=L^{-(d-1)/2}$、$[r]=L^{-2}$、$[u]=L^{d-3}$。粗粒化长度放大$b$时无量纲四次耦合按$u'\simeq b^{3-d}u$变换：$d<3$相关，$d=3$边缘，$d>3$无关。这里的上临界\emph{空间}维数是$d_c=3$，因为量子时间还贡献一维，时空维数$D=d+1=4$；经典静态$\phi^4$模型的上临界空间维数才是4。$d=1$链已知的精确Majorana解不能被高维Gaussian幂次计数替代。

\originalsection{环上的边界与宇称投影}
周期自旋链的最后一条键穿过Jordan--Wigner串，费米边界条件取决于总宇称$P=\prod_j\sigma_j^z=(-1)^{N_f}$。以上约定通常得到$c_{L+1}=-Pc_1$：偶宇称采用反周期，奇宇称采用周期。每个费米边界模型还须投影到匹配宇称，不能把两个完整Fock空间相加，否则会把原$2^L$个自旋态双计。

反周期允许$k=2\pi(n+1/2)/L$，所有$k$可两两配对，Bogoliubov真空偶宇称。周期有$k=0,\pi$自配对点，其占据由$h-J$与$h+J$的符号决定，在跨越临界点时可改变真空宇称。非匹配扇区的最低物理态要额外加入准粒子。有限环的最低能态选择因此不能只对两个动量格点求和后忽略宇称约束。

有序区两种自旋宇称最低态的能量差为指数小而非严格零。能谱$\sqrt{J^2+h^2-2Jh\cos k}$在复$k$平面的最近支点满足$\cosh\kappa=(J^2+h^2)/(2Jh)$，故$\kappa=|\log(J/h)|$。Poisson求和把两种边界格点之差写成高Fourier系数，最近支点给$\e^{-\kappa L}$与可能的代数前因子，这与端模局域长度一致。临界点支点到实轴，指数估计失效，有限尺寸能隙改为$1/L$。
\originalsection{对偶与局域可观测量}
定义对偶键自旋$\mu_{j+1/2}^z=\sigma_j^x\sigma_{j+1}^x$、$\sigma_j^z=\mu_{j-1/2}^x\mu_{j+1/2}^x$，交换$J,h$后哈密顿量形状保持，临界点在自对偶处。环上全局约束与边界仍必须跟踪，对偶不在无条件下把任意边界Hilbert空间一一对应。

自旋序参数$\sigma^x$包含费米字符串；$\sigma^z$是局域费米双线性。因而动态$\sigma^z$相关通常可创建两费米子，谱连续边缘由$E_k+E_{q-k}$决定。若探针改变自旋宇称，需把它作用后的边界扇区一并处理。直接把$E_k$画成中子散射的唯一$\delta$峰会漏掉算符选择规则。
